\chapter{KẾT LUẬN VÀ ĐỀ XUẤT}\label{chap4_intro}
\vspace{-40pt}

\section{Kết luận và mức độ hoàn thành đề tài}

\subsection{Tổng hợp kết quả đạt được so với mục tiêu}

Về phân tích và thiết kế, báo cáo đã xác định các tác nhân gồm khách vãng lai, khách hàng, nhân viên, quản trị viên, cổng thanh toán và dịch vụ Google Identity; đồng thời mô tả 18 ca sử dụng cho các nghiệp vụ từ tra cứu sản phẩm đến quản lý voucher. Sơ đồ Use Case tổng quan, các sơ đồ trình tự và sơ đồ cơ sở dữ liệu giúp liên kết yêu cầu với các bước xử lý và dữ liệu chính như tài khoản, sản phẩm, giỏ hàng, đơn hàng, địa chỉ, tồn kho, khuyến mãi và giao dịch thanh toán. Đây là cơ sở phù hợp để tổ chức các chức năng trong phạm vi website bán máy tính và linh kiện.

Về xây dựng hệ thống, mã nguồn dùng PHP, MySQL và PDO, với giao diện HTML, CSS và JavaScript; báo cáo xác định môi trường chạy là Apache trong XAMPP. Front controller và router chuyển yêu cầu đến controller, các model xử lý dữ liệu, còn view hiển thị nội dung. Cách tổ chức này đã kết nối các trang dành cho khách hàng với nghiệp vụ và cơ sở dữ liệu, đồng thời tách khu vực quản trị khỏi trải nghiệm mua sắm. Giao diện có các trang danh mục, tìm kiếm, chi tiết sản phẩm, giỏ hàng, thanh toán, tài khoản và quản trị; mã CSS cũng có các quy tắc hiển thị theo kích thước màn hình.

Ở phía khách hàng, hệ thống đã triển khai đăng ký, đăng nhập bằng tài khoản website, đăng nhập hoặc liên kết tài khoản Google, cập nhật hồ sơ và sổ địa chỉ. Khách có thể tìm theo tên, lọc sản phẩm theo danh mục, thương hiệu, loại và khoảng giá, sắp xếp danh sách, xem chi tiết, quản lý giỏ hàng và áp dụng khuyến mãi hoặc voucher khi đặt hàng. Đơn hàng lưu thông tin nhận hàng và giá trị giảm tại thời điểm tạo; khách hàng có thể xem trạng thái và yêu cầu hủy khi đơn còn ở trạng thái Chờ xác nhận. Các bước kiểm tra giá, voucher và tồn kho được thực hiện lại ở phía máy chủ.

Ở phía vận hành, nhân viên và quản trị viên có các chức năng phù hợp với vai trò để xử lý đơn và xem tồn kho; quản trị viên còn quản lý sản phẩm, danh mục, người dùng, nhân viên, kho, chương trình giảm giá, voucher và thống kê. Tồn kho được cập nhật cùng với việc tạo, hủy hoặc xử lý hàng hoàn của đơn, và lịch sử kho ghi nhận sản phẩm, lượng thay đổi, thao tác, người thực hiện và thời gian. Thống kê hiện tập trung vào số đơn theo trạng thái, doanh thu từ đơn hoàn thành, sản phẩm bán nhiều và sản phẩm sắp hết hàng; chưa có cơ sở để mô tả đây là báo cáo lợi nhuận.

Một số biện pháp bảo vệ dữ liệu đã được triển khai trong mã nguồn: truy vấn tham số hóa bằng PDO, băm và kiểm tra mật khẩu, kiểm tra CSRF, xác thực phiên và vai trò ở controller, cùng transaction và khóa bản ghi cho một số nghiệp vụ như đặt hàng, cập nhật tồn kho, voucher và ghi nhận thanh toán. Luồng trực tuyến qua SePay có phần khởi tạo checkout và tiếp nhận IPN; mã xử lý kiểm tra thông tin xác thực, trạng thái, tiền tệ và số tiền trước khi cập nhật đơn. Tuy nhiên, đây là tích hợp Sandbox; việc có mã xử lý không tự chứng minh giao dịch đã được kiểm thử xuyên suốt với cấu hình thật hoặc sẵn sàng vận hành production.

Về kiểm thử, Chương 3 hiện tập trung vào đặc tả chức năng và chưa có bảng ca kiểm thử, nhật ký hoặc kết quả thực tế để đối chiếu. Vì vậy, không thể xác nhận tổng số ca đạt, không đạt hay chưa thực hiện, cũng không thể tính tỷ lệ đạt từ tài liệu hiện có. Chưa có căn cứ để kết luận về hiệu năng, an toàn trước tấn công, khả năng chịu tải hay độ ổn định khi vận hành dài hạn.

\subsection{Phạm vi đã hoàn thành và những nội dung còn giới hạn}

Trong phạm vi Đồ án 1, các chức năng cốt lõi đã hình thành một quy trình mua hàng có thể theo dõi từ tra cứu sản phẩm, quản lý giỏ và thông tin nhận hàng đến tạo đơn; phía quản trị có thể quản lý dữ liệu bán hàng, điều chỉnh tồn kho, xử lý trạng thái đơn và xem các thống kê cơ bản. Việc tính lại giá và voucher ở máy chủ, cập nhật tồn kho cùng đơn hàng trong transaction, và ghi lịch sử thay đổi giúp giảm nguy cơ dữ liệu đơn lẻ bị lệch khi một thao tác thất bại. Đây là mức hoàn thành phù hợp cho một hệ thống đồ án chạy trong môi trường phát triển.

Một số nội dung được mô tả trong thiết kế chưa thấy triển khai đầy đủ. Chương 2 nêu hệ thống tự hủy đơn khi thanh toán thất bại hoặc quá thời hạn, nhưng trong các route và mã nguồn đã rà soát chưa thấy tác vụ nền hoặc lịch chạy để thực hiện việc tự hủy; hiện có luồng khách hàng tự hủy khi đơn còn chờ xác nhận và luồng thanh toán trực tuyến thất bại giữ đơn ở trạng thái chưa thanh toán để thử lại. Vì vậy, quy tắc tự hủy theo thời hạn cần được xem là yêu cầu chưa hoàn tất, không phải chức năng đã xác nhận hoạt động.

Các trang hướng dẫn bảo hành, đổi trả và giao hàng hiện cung cấp nội dung thông tin; chưa thấy quy trình tạo, phân công và theo dõi yêu cầu bảo hành hoặc đổi trả của khách hàng. Trạng thái giao hàng được cập nhật trong quản trị, nhưng chưa thấy tích hợp hãng vận chuyển để đồng bộ mã vận đơn và hành trình. Tương tự, hệ thống có thống kê doanh thu, đơn hàng, sản phẩm bán chạy và tồn thấp nhưng chưa thấy mô hình dữ liệu giá vốn đủ để tính lợi nhuận. Những điểm này là giới hạn chức năng hiện tại, không nên diễn giải thành các nghiệp vụ đã hoàn thành.

Một giới hạn khác nằm ở mức xác nhận kết quả. Báo cáo chưa cung cấp nhật ký kiểm thử có thể đối chiếu, và chưa ghi nhận chạy kiểm thử tích hợp trên Apache/MySQL với dữ liệu thực. Google Identity và SePay phụ thuộc cấu hình, dịch vụ ngoài và điều kiện mạng; mã nguồn thể hiện các điểm tích hợp nhưng chưa có bằng chứng chạy thử đầy đủ trong tài liệu. Mã CSS có các breakpoint cho màn hình nhỏ, song chưa có ma trận kiểm tra trình duyệt và thiết bị để xác nhận trải nghiệm đồng nhất.

Tài liệu cũng cần thống nhất phạm vi thanh toán giữa các chương: Chương 1 nêu COD là phương thức được triển khai, trong khi Chương 2 mô tả thanh toán trực tuyến và mã nguồn có luồng SePay Sandbox. Kết luận hợp lý là chức năng COD có trong luồng đặt hàng, còn thanh toán trực tuyến mới được tích hợp ở môi trường thử nghiệm và chưa có bằng chứng nghiệm thu production. Nhìn chung, nhóm đã hiện thực hóa phần lớn luồng mua hàng và quản trị cơ bản, nhưng mức độ hoàn thành chưa bao gồm đầy đủ quy trình sau bán hàng, tác vụ tự hủy đơn theo hạn và kiểm chứng vận hành thực tế.

\section{Hướng phát triển và kiến nghị}

\subsection{Hoàn thiện chức năng nghiệp vụ và trải nghiệm mua sắm}

Danh mục hiện hỗ trợ tìm theo tên, lọc danh mục, thương hiệu, loại sản phẩm và khoảng giá, nhưng chưa có so sánh cấu hình giữa các mẫu hoặc tìm kiếm toàn văn trên thông số. Với máy tính và linh kiện, khách hàng thường cần đối chiếu CPU, GPU, bộ nhớ và khả năng tương thích; vì vậy có thể bổ sung bộ lọc theo thuộc tính kỹ thuật đã chuẩn hóa và chức năng so sánh một số sản phẩm trên cùng bảng. Cách này giúp khách đánh giá lựa chọn dựa trên thông số thay vì chỉ dựa vào tên và giá.

Trang chi tiết đã cung cấp thông tin sản phẩm, nhưng mức hữu ích của trang phụ thuộc vào dữ liệu kỹ thuật được nhập đầy đủ và đồng nhất. Nhóm có thể xác định bộ thuộc tính theo từng loại sản phẩm, kiểm tra trường bắt buộc khi quản trị viên lưu sản phẩm, và trình bày thông số theo cấu trúc thống nhất. Sau đó, đánh giá hoặc nhận xét chỉ nên mở khi có cơ chế gắn với đơn đã mua, kiểm duyệt nội dung và xử lý báo cáo vi phạm; điều này giúp phản hồi có ngữ cảnh và hạn chế đánh giá không đáng tin cậy.

Đơn hàng hiện được theo dõi bằng trạng thái nội bộ, còn thông tin bảo hành, đổi trả và giao hàng chủ yếu là các trang hướng dẫn. Hướng phát triển phù hợp là gắn từng yêu cầu sau bán hàng với mã đơn hoặc sản phẩm cụ thể, lưu lý do, người xử lý, thời điểm và kết quả; đồng thời bổ sung lịch sử trạng thái để khách biết yêu cầu đang ở bước nào. Khi có đơn vị vận chuyển phù hợp, có thể tích hợp mã vận đơn và cập nhật trạng thái giao nhận. Các giao diện hiện đã có CSS theo kích thước màn hình, nên cần kiểm tra thực tế các trang danh mục, giỏ và quản trị trên điện thoại, máy tính bảng và một số trình duyệt trước khi khẳng định khả năng tương thích.

\subsection{Mở rộng quản lý và tích hợp dịch vụ}

Trang thống kê hiện cung cấp số đơn, trạng thái đơn, doanh thu của đơn hoàn thành, sản phẩm bán chạy và danh sách tồn thấp. Để hỗ trợ quyết định nhập hàng và đánh giá hoạt động tốt hơn, có thể bổ sung xuất báo cáo theo khoảng thời gian và đối chiếu các chỉ tiêu với dữ liệu đơn hàng. Nếu cần tính lợi nhuận, trước hết phải xác định và lưu giá vốn đáng tin cậy theo từng lần nhập hoặc sản phẩm; chỉ lấy doanh thu trừ một con số ước tính sẽ tạo ra kết quả không đủ căn cứ.

Quản trị kho đã có thao tác điều chỉnh và lịch sử thay đổi, nhưng chưa thấy chức năng kiểm kê để so sánh số lượng thực tế với số lượng ghi nhận. Có thể bổ sung phiên kiểm kê theo sản phẩm, lưu người kiểm, thời điểm, chênh lệch và lý do điều chỉnh; báo cáo đối chiếu sẽ giúp phát hiện sai lệch và truy nguyên thay đổi. Cùng với đó, cần hoàn thiện quy trình nhập liệu hoặc nhập danh sách sản phẩm có kiểm tra mã, danh mục và thuộc tính để giảm lỗi do thao tác thủ công.

SePay hiện được cấu hình theo Sandbox. Trước khi cân nhắc sử dụng môi trường thật, cần kiểm thử các trường hợp callback/IPN lặp, sai chữ ký, sai số tiền, giao dịch đến muộn và đơn đã bị hủy; đồng thời đối chiếu bản ghi giao dịch với trạng thái đơn và có quy trình xử lý giao dịch chưa rõ kết quả. Việc ghi nhận giao dịch đã có trong schema và mã nguồn, nhưng vận hành thật còn cần cấu hình chính thức, giám sát lỗi và quy trình đối soát. Các bước này giúp hạn chế việc ghi nhận nhầm hoặc bỏ sót thanh toán.

\subsection{Nâng cao bảo mật, hiệu năng và khả năng vận hành}

Mã nguồn đã có prepared statements, CSRF, băm mật khẩu và kiểm tra vai trò ở phía máy chủ; các biện pháp này tạo nền tảng nhưng chưa thay thế một đợt rà soát bảo mật. Nên lập danh sách quyền theo từng route và vai trò, kiểm tra truy cập trực tiếp khi thiếu phiên hoặc sai quyền, rà soát việc mã hóa đầu ra tại view và cấu hình cookie phiên. Các khóa dịch vụ cần được lưu trong cấu hình triển khai riêng, không đưa vào kho mã nguồn; nhật ký cần ghi đủ sự kiện quan trọng mà không chứa thông tin xác thực nhạy cảm.

Truy vấn danh mục hiện dùng điều kiện tìm theo tên, kết hợp bộ lọc và phép tính giá hiệu lực của khuyến mãi. Khi dữ liệu tăng, nhóm nên đo thời gian truy vấn trên dữ liệu có quy mô đại diện, xem kế hoạch truy vấn và chỉ mục cho các cột lọc, sau đó tối ưu các phép tìm kiếm hoặc tính giá có chi phí cao. Việc kiểm thử tải nên bao gồm nhiều khách đồng thời đặt cùng sản phẩm, cập nhật tồn kho và gửi lại IPN; tiêu chí quan trọng là không bán vượt tồn, không tạo giao dịch thanh toán trùng và không làm sai trạng thái đơn.

Trước khi triển khai ngoài môi trường phát triển, cần dựng một môi trường thử nghiệm gần với máy chủ dự kiến, kiểm tra cấu hình PHP, Apache, MySQL, HTTPS và các tích hợp ngoài, rồi mới chuyển sang vận hành có giới hạn. Nhóm cũng nên thiết lập sao lưu cơ sở dữ liệu định kỳ và thực hành khôi phục trên bản sao để kiểm tra khả năng phục hồi, thay vì chỉ tạo tệp sao lưu mà chưa kiểm chứng. Cuối cùng, cần xây dựng bộ ca kiểm thử có dữ liệu đầu vào, kết quả mong đợi, kết quả thực tế và bằng chứng cho đăng nhập, phân quyền, đặt/hủy đơn, tồn kho, voucher, thanh toán và lỗi dịch vụ; khi đó mới có thể báo cáo tỷ lệ đạt một cách minh bạch.

\addtocontents{toc}{\protect\setcounter{tocdepth}{0}}
